\documentclass[11pt]{article}
\usepackage{amsmath,amssymb,amsthm}
\usepackage{hyperref}
\usepackage{geometry}
\geometry{margin=1in}

\newtheorem{theorem}{Theorem}
\newtheorem{lemma}[theorem]{Lemma}
\newtheorem{proposition}[theorem]{Proposition}

\newcommand{\GF}{\mathbb{F}_2}
\newcommand{\ket}[1]{\left|#1\right\rangle}

\title{Exact Entanglement Entropy and Code Distance for Stabilizer States,\\
Computed Without Floating-Point Arithmetic}
\author{Joel Cruz Cabrera \\ \href{https://orcid.org/0009-0005-4048-1237}{ORCID: 0009-0005-4048-1237}}
\date{2026}

\begin{document}
\maketitle

\begin{abstract}
Standard quantum circuit simulation represents amplitudes as
double-precision complex numbers, so quantities such as entanglement
entropy are computed only approximately. We revisit two well-known facts
about the stabilizer formalism --- that entanglement entropy of a
stabilizer state is always an integer, and that the distance of a
stabilizer code is a combinatorial quantity --- and implement both
\emph{exactly}, using binary linear algebra over $\GF$ on the stabilizer
tableau, with no floating-point step anywhere in the computation. We
verify the entropy formula against Qiskit's floating-point von Neumann
entropy on canonical states and on 30 random Clifford circuits (all
matches exact), reproduce the area-law entanglement of the 1D cluster
state, and verify the code-distance procedure against three canonical
quantum error-correcting codes (Steane, the 5-qubit code, and Shor's
code), recovering their published distances exactly. As an illustration
of the exact toolkit in use, we report an empirical scaling law for the
depth at which entanglement saturates in a specific random-circuit
model. Finally, since every result above concerns exact mathematics or
exact simulation of noiseless Clifford circuits, we complement it with a
small, honestly-labelled real-hardware comparison: the GHZ state on
$3$--$7$ qubits, run on IBM Quantum superconducting hardware, measured
against the exact noiseless prediction our formalism computes, together
with a quantitative check against a naive noise model built from the
device's own published calibration data --- which the data
systematically contradicts by 18--21 percentage points, a discrepancy we
report along with our best physical explanation for it. A second,
independent hardware experiment measures the cluster state's own
stabilizer generators directly, and identifies a single miscalibrated
physical qubit as correlated with elevated error; a follow-up control
experiment, forcing the circuit onto verified low-error qubits instead,
confirms the effect is causal but partial (leakage drops by $1.84$
percentage points, roughly a third of the excess) rather than the whole
explanation, which we report at its true, unexaggerated size. Applying
the naive noise model to the control circuit itself reverses its bias
(now optimistic rather than pessimistic) by an amount consistent, to
order of magnitude, with ordinary decoherence over the circuit's
execution time --- a physical effect the naive model was never built to
capture.
\end{abstract}

\noindent \textbf{MSC codes:} 81P45, 81P70, 94B05, 68Q12.

\noindent \textbf{Keywords:} stabilizer formalism, Gottesman--Knill
theorem, entanglement entropy, quantum error-correcting codes, exact
computation.

\section{Introduction}

A generic $n$-qubit quantum state is described by $2^n$ complex
amplitudes, and quantities derived from it --- entanglement entropy in
particular --- are irrational in general and are computed numerically,
with floating-point rounding error. The Gottesman--Knill theorem
identifies a restricted but important class of states and circuits,
those generated by the Clifford group (Hadamard, phase, and CNOT gates),
for which the state admits an \emph{exact}, discrete representation: an
$n \times 2n$ binary matrix of stabilizer generators, with no continuous
amplitudes at all \cite{Gottesman1997,NielsenChuang}. This is the
mathematical content behind the classical simulability of Clifford
circuits, and it underlies the entire theory of stabilizer quantum
error-correcting codes.

Two quantities that are usually computed numerically for general
states admit exact combinatorial descriptions in the stabilizer case:

\begin{itemize}
\item The \emph{entanglement entropy} of a stabilizer state across any
  bipartition is an integer multiple of $\log 2$ --- never irrational
  --- computable as a rank over $\GF$ of a submatrix of the stabilizer
  tableau \cite{Fattal2004}.
\item The \emph{distance} of a stabilizer code is the minimum weight of
  a nontrivial logical Pauli operator, computable by an exhaustive
  search over the Pauli group restricted to the code's centralizer
  \cite{Gottesman1997}.
\end{itemize}

Neither fact is new; both are standard in the quantum information
literature. What we report here is a from-scratch, dependency-light
implementation of both quantities using pure binary linear algebra ---
no numerical linear algebra library, no floating-point arithmetic at any
point --- together with an explicit verification protocol: every claim
below was checked against an independent computation (Qiskit's
floating-point statevector simulator, or a published code's known
parameters) before being trusted.

\section{Exact entanglement entropy}

Let $\ket{\psi}$ be an $n$-qubit stabilizer state with independent
stabilizer generators $g_1, \ldots, g_n$, each written in symplectic
form as a pair $(x_i, z_i) \in \GF^n \times \GF^n$ (so $g_i$ acts as $X$
on qubit $q$ if $x_{i,q}=1$, as $Z$ if $z_{i,q}=1$, as $Y$ if both, and
trivially otherwise).

\begin{theorem}[Fattal, Cubitt, Yamamoto, Bravyi, Chuang, 2004
\cite{Fattal2004}]
For a bipartition into a region $A$ (with complement $B$), the
entanglement entropy of $\ket{\psi}$ across $A|B$, in units of $\log 2$,
is
\[
  S_A \;=\; \operatorname{rank}_{\GF}(M_A) - |A|,
\]
where $M_A$ is the $n \times 2|A|$ matrix formed by taking, from each
generator, only the $x$ and $z$ components restricted to the qubits in
$A$.
\end{theorem}

We implement $\operatorname{rank}_{\GF}$ by Gaussian elimination with
XOR row operations on Python integers/booleans, avoiding NumPy or any
floating-point library. The stabilizer tableau itself is read directly
from Qiskit's \texttt{Clifford} object (\texttt{stab\_x}, \texttt{stab\_z}),
which stores it as a boolean array --- exact by construction, since a
Clifford circuit's action on the Pauli group is itself a permutation
with signs, requiring no continuous arithmetic.

\subsection{Verification}

We checked $S_A$ against Qiskit's von Neumann entropy
(\texttt{Statevector} + \texttt{partial\_trace} + \texttt{entropy}, a
floating-point computation, used here purely as an independent oracle,
never as the reported value) on:

\begin{itemize}
\item The product state $\ket{00}$: $S_{\{0\}} = 0$.
\item The Bell pair $(\ket{00}+\ket{11})/\sqrt2$: $S_{\{0\}} = 1$.
\item The GHZ state on $n$ qubits, cut at a single qubit: $S_{\{0\}} =
  1$ for $n = 2,3,5,8$.
\item The one-dimensional cluster (graph) state on $n$ qubits: $S_A = 1$
  for \emph{every} bulk cut, for $n = 4, 6, 8$ --- the area-law
  signature expected of this state.
\item Thirty random Clifford circuits ($n=5$ qubits, 20 random gates
  each, uniformly random bipartition): all 30 exact values matched the
  floating-point reference to the precision of the float computation.
\end{itemize}

\section{Exact code distance}

For a stabilizer code with $n-k$ independent generators (so $k$ logical
qubits), Qiskit's \texttt{Clifford} object is not directly applicable,
since it represents a full $n$-qubit Clifford unitary rather than a code
with $k>0$. We instead work directly with the symplectic representation.

\begin{proposition}
The distance $d$ of a stabilizer code is the minimum Hamming weight,
over qubits, of a Pauli operator $P$ satisfying:
\begin{enumerate}
\item $P$ commutes with every stabilizer generator (symplectic inner
  product zero against each), and
\item $P$ is not itself a product of stabilizer generators (i.e., its
  symplectic vector does not lie in the $\GF$-span of the generators'
  vectors).
\end{enumerate}
\end{proposition}

We search candidates by increasing weight $w = 1, 2, \ldots$, enumerating
all $\binom{n}{w} 3^w$ Pauli strings of that weight, and return the first
$w$ for which a valid logical operator is found. Membership in the span
is tested by reduction against a row-echelon basis of the generator
matrix, again over $\GF$ with no floating point.

\subsection{Verification against known codes}

We implemented three canonical codes independently (not imported from
any coding-theory library) and computed their distance:

\begin{center}
\begin{tabular}{lccc}
Code & $[[n,k,d]]$ published & $d$ computed & Match \\
\hline
Steane (CSS on Hamming $[7,4,3]$) & $[[7,1,3]]$ & 3 & Yes \\
5-qubit perfect code & $[[5,1,3]]$ & 3 & Yes \\
Shor's code & $[[9,1,3]]$ & 3 & Yes \\
\end{tabular}
\end{center}

These are known results; the point of the check is that our
independent, from-scratch implementation reproduces them exactly, which
we take as validation of the method for use on codes whose distance is
not already tabulated.

\section{An illustration: entanglement saturation depth}

As a demonstration of the exact toolkit applied to a genuine question,
we measured how quickly entanglement grows under a specific random
circuit model: at each step, a gate is either a random single-qubit
Clifford gate (probability $1/2$) or a CNOT between a uniformly random
pair of qubits, drawn from \emph{all} $\binom{n}{2}$ pairs, not only
nearest neighbors (probability $1/2$). We tracked $S_A(t)$, exactly, at
every step, for a half--half bipartition, and recorded the first step
at which $S_A(t)$ reaches its maximum possible value $n/2$.

Averaging over 10 random trajectories at each of $n = 10, 20, 40, 80$
qubits, the mean saturation depth $t_{\text{sat}}$ and the ratio
$t_{\text{sat}}/n$ were:

\begin{center}
\begin{tabular}{ccc}
$n$ & mean $t_{\text{sat}}$ & $t_{\text{sat}}/n$ \\
\hline
10 & 144.4 & 14.44 \\
20 & 293.9 & 14.69 \\
40 & 553.5 & 13.84 \\
80 & 1130.1 & 14.13 \\
\end{tabular}
\end{center}

The ratio is stable across a factor-of-8 range in $n$ (within
$\pm 0.5$ of 14.3), consistent with linear scaling $t_{\text{sat}}(n)
\approx 14.3\, n$ for this specific circuit model. A plausible heuristic:
with uniform pairing over all qubits, the probability that a given step
crosses a half--half cut is $\Theta(1)$, independent of $n$, so the
$\Theta(n)$ bits of entanglement needed to saturate should require
$\Theta(n)$ steps. We report this as a numerical observation with a
plausible but unproven heuristic explanation, not as a theorem; four
data points with ten trajectories each do not rule out other functional
forms (e.g., $n\log n$), and the specific constant is a property of this
circuit model, not a universal quantity.

\section{A real-hardware comparison}\label{sec:hardware}

Everything above is either exact mathematics or exact simulation of
noiseless Clifford circuits; no physical noise enters at any point. To
put a number on how far actual current hardware is from the states this
formalism describes, we ran the simplest circuit in the paper --- the
GHZ state, $\ket{\text{GHZ}(n)} = (\ket{0^n}+\ket{1^n})/\sqrt2$, prepared
by \texttt{ghz\_state(n)} from Section~2 followed by a computational-basis
measurement --- on real hardware, for $n=3,\ldots,7$.

\subsection{Exact prediction and setup}

For a noiseless GHZ state, a computational-basis measurement gives
$P(0^n) = P(1^n) = 1/2$ and $P(\text{anything else}) = 0$ exactly, by
the definition of the state; this requires no computation beyond reading
off the amplitudes. Circuits were transpiled and executed via IBM Quantum
Runtime (\texttt{SamplerV2}) on \texttt{ibm\_kingston}, a $156$-qubit
superconducting processor, at $4096$ shots per circuit.

\subsection{Results}

\begin{center}
\begin{tabular}{lccc}
Circuit & Job ID (IBM Quantum) & $P(0^n)+P(1^n)$ measured & Exact \\
\hline
GHZ(3), run 1 & \texttt{d9kuc78ii2cc73egnv4g} & 0.9441 & 1.0000 \\
GHZ(3), run 2 & \texttt{d9kukujhdfks73ck7g0g} & 0.9412 & 1.0000 \\
GHZ(3), run 3 & \texttt{d9kul9bjf64c739iolog} & 0.9402 & 1.0000 \\
GHZ(4) & \texttt{d9kumsqbr2fc73e7tsvg} & 0.9307 & 1.0000 \\
GHZ(5) & \texttt{d9kuolabr2fc73e7tvf0} & 0.9167 & 1.0000 \\
GHZ(6) & \texttt{d9kv0tibr2fc73e7udjg} & 0.8896 & 1.0000 \\
GHZ(7) & \texttt{d9kv4ojjf64c739ipdgg} & 0.8828 & 1.0000 \\
\end{tabular}
\end{center}

The three GHZ(3) repetitions agree closely with one another (leakage to
non-ideal outcomes: $5.59\%$, $5.88\%$, $5.98\%$; mean $5.82\%$, sample
standard deviation $0.20\%$), indicating the hardware's error rate on
this specific circuit was stable across separate job submissions, not
merely noisy from run to run. Across $n = 3,\ldots,7$, the leakage
increases monotonically ($5.8\% \to 6.9\% \to 8.3\% \to 11.0\% \to
11.7\%$), consistent with the expectation that a longer CNOT chain and
more measured qubits accumulate more gate and readout error --- but we
report this as a qualitative, seven-data-point observation, not a fitted
noise model.

\subsection{A quantitative check: does a naive error model predict this?}

The qualitative trend above invites an obvious quantitative question:
does simply multiplying IBM's own published per-gate and per-measurement
error rates, for the exact physical qubits each transpiled circuit
actually used, predict the observed leakage? We pulled
\texttt{ibm\_kingston}'s calibration data (via \texttt{backend.target})
on the same day the circuits ran, and computed, for each job,
\[
  P_{\text{naive}} \;=\; \prod_{\text{gate/measurement } g \text{ used}}
  \bigl(1 - \varepsilon_g\bigr),
\]
where each $\varepsilon_g$ is IBM's own published error rate for that
specific gate on those specific physical qubits (a standard first-order
circuit-fidelity estimate, not a fitted model).

\begin{center}
\begin{tabular}{lccc}
$n$ & Job ID & $P_{\text{naive}}$ predicted & $P(0^n)+P(1^n)$ measured \\
\hline
3 (run 1) & \texttt{d9kuc78ii2cc73egnv4g} & 0.7337 & 0.9441 \\
3 (run 2) & \texttt{d9kukujhdfks73ck7g0g} & 0.7337 & 0.9412 \\
3 (run 3) & \texttt{d9kul9bjf64c739iolog} & 0.7337 & 0.9402 \\
4 & \texttt{d9kumsqbr2fc73e7tsvg} & 0.7286 & 0.9307 \\
5 & \texttt{d9kuolabr2fc73e7tvf0} & 0.7187 & 0.9167 \\
6 & \texttt{d9kv0tibr2fc73e7udjg} & 0.7104 & 0.8896 \\
7 & \texttt{d9kv4ojjf64c739ipdgg} & 0.7025 & 0.8828 \\
\end{tabular}
\end{center}

The naive model \emph{underestimates} survival in all seven runs, by a
systematic $18$--$21$ percentage points --- not scattered error, the same
direction and a similar size every time. We checked, and ruled out, the
easy explanation: the Runtime job options were submitted empty, and
\texttt{SamplerV2}'s twirling and dynamical-decoupling options default to
disabled, so no hidden error mitigation was applied. The more likely
explanation is a known property of the underlying calibration data
itself: randomized-benchmarking error rates measure a gate's total
average infidelity against the ideal unitary, over \emph{all} error
types (including phase errors), not specifically the probability that a
given error flips the bit ultimately read out in a computational-basis
measurement. Since GHZ preparation is dominated by two-qubit gates whose
errors are often phase-type, a naive bit-flip-equivalent model built
from these published rates is expected to be pessimistic --- which is
exactly the direction of the discrepancy observed. We report this as a
plausible explanation consistent with known physics, not as something we
verified by a separate, independent measurement of phase- versus
bit-flip-error content.

One calibration detail is worth recording on its own: physical qubit
$1$, used in all three GHZ(3) runs, had a published readout error of
$23.4\%$ that day --- against a median of $0.9\%$ across
\texttt{ibm\_kingston}'s $156$ qubits (eight qubits exceeded $10\%$).
Despite including this outlier qubit, the three GHZ(3) runs still showed
only $5.6$--$6.0\%$ total leakage, a further concrete illustration that a
single large published error rate does not translate one-to-one into
observed bit-flip probability in this kind of experiment.

\subsection{A causal test: avoiding the miscalibrated qubit}

The observation above is correlational: qubit $1$ was present, and the
leakage was elevated. To test the claim directly rather than merely
observe it, we repeated the GHZ(3) circuit with the transpiler's
\texttt{initial\_layout} forced onto a different, independently verified
low-error chain of physical qubits ($118$--$129$--$128$; readout errors
$0.29\%$, $0.34\%$, $0.42\%$, all well below the chip median), explicitly
avoiding qubit $1$. We confirmed the transpiler respected the requested
layout before submitting.

\[
  \text{control (no qubit }1\text{): leakage } 3.98\% \qquad\text{vs.}\qquad
  \text{baseline (qubit }1\text{ included): mean } 5.82\%,\ \text{s.d. } 0.20\%.
\]

The improvement ($1.84$ percentage points) is real --- roughly nine
times the baseline's own run-to-run standard deviation, not noise --- but
it is \emph{partial}, not the near-total elimination of leakage that
would follow if the single qubit explained everything. Avoiding it
accounts for roughly a third of the original leakage; a residual
$\sim\!4\%$ leakage floor remains even on the best available qubits. We
report the honest magnitude of the effect, not the more dramatic story a
smaller number would have made: qubit miscalibration is a real, causally
confirmed contributor, and not the whole explanation.

\subsection{Where does the residual floor come from?}

Applying the same naive error-rate-multiplication model of
Section~5.3 to the control circuit itself (physical qubits $118,129,128$)
gives a striking reversal: the naive model predicts $98.53\%$ survival
($1.47\%$ leakage) --- but the measured leakage was $3.98\%$. For the
qubit-$1$ runs the naive model was too \emph{pessimistic}; here, on
good-quality qubits, it is too \emph{optimistic}. The direction of the
naive model's error flips with qubit quality, which is itself informative:
the model has a blind spot that does not depend on which qubits are used.

That blind spot is idle-time decoherence. The naive model multiplies only
\emph{per-operation} error rates; it never accounts for the fact that a
qubit must remain coherent for the circuit's entire duration, not just
during the instants a gate acts on it. A rough estimate: summing the
published durations of every gate and measurement in the actually-submitted
control circuit gives $6.84\,\mu\text{s}$ (treated as fully serial, an
overestimate relative to real parallelism); dividing by the shortest
$T_1$ among the three qubits ($T_1 = 230\,\mu\text{s}$ for physical qubit
$128$) gives a decoherence-probability estimate of $2.97\%$ --- close to
the $2.51$-point gap between the naive prediction and what was measured.
This is a crude estimate, not a rigorous channel model (real parallelism
would shorten the effective duration; $T_2$, not $T_1$, may be the more
relevant timescale for a phase-sensitive entangled state, and physical
qubit $128$'s published $T_2 = 63.3\,\mu\text{s}$ is markedly shorter
than its $T_1$), but the order-of-magnitude agreement is consistent with
a clear conclusion: the residual floor looks like ordinary decoherence
during circuit execution, a physical effect the naive per-gate model was
never built to capture, rather than an unexplained anomaly.

\subsection{What this does and does not show}

This is not a claim of quantum advantage, nor a new theorem, nor a
characterisation of the device's noise in general: it is a single dated
snapshot (job IDs above, executed 2026-07-29) of one specific device
running one specific shallow circuit. Reproducing these exact percentages
on a different day, or on different hardware, is not expected; the
qubit-count trend is offered only as a plausible, physically unsurprising
pattern, not a law. The control experiment is the one causal claim in
this section, and even there we report a partial, bounded effect rather
than overstate it as the full explanation. Likewise, the naive-model
discrepancy above is a
single, consistent seven-point observation with a physically plausible
explanation, not a validated noise model of the device --- someone with
access to the device's phase-error and bit-flip-error content separately
could test the explanation directly, which we did not do here. What the
experiment does establish is procedural:
the same circuit-construction code already verified exactly in
Section~2 runs unmodified on real hardware, and the exact, noiseless
prediction this paper computes throughout is a meaningful, well-defined
baseline against which real devices can be honestly measured.

\section{Cluster-state stabilizers on real hardware}\label{sec:cluster-hardware}

The GHZ comparison above tests only whether a computational-basis
measurement lands on one of two ideal bitstrings. The cluster state
studied in Section~4 offers a more direct test of the stabilizer
structure this paper is built on: for the linear cluster state on qubits
$0,\ldots,n-1$ (prepared by \texttt{cluster\_state\_1d(n)}), each qubit
$i$ has an associated stabilizer generator
\[
  K_0 = X_0 Z_1, \qquad
  K_i = Z_{i-1} X_i Z_{i+1} \ (0<i<n-1), \qquad
  K_{n-1} = Z_{n-2} X_{n-1},
\]
and, by the definition of a stabilizer state, $\langle K_i \rangle = 1$
\emph{exactly} for every $i$. We measured all $n=4$ generators on real
hardware: rotate qubit $i$ into the $Z$ basis (an $H$ gate) before
measuring every qubit in the computational basis, then estimate
$\langle K_i\rangle$ as the shot-averaged product of $(-1)^{\text{bit}}$
over the qubits in $K_i$'s support. Before spending real hardware time,
we checked this circuit-and-analysis pipeline against a noiseless exact
simulation (Qiskit's \texttt{Statevector}, both the analytic expectation
value and shot-sampled counts): all four generators reproduced
$\langle K_i\rangle = 1.0000$ exactly, confirming the measurement logic
before it was trusted on hardware.

\begin{center}
\begin{tabular}{lccc}
Generator & Support & $\langle K_i\rangle$ measured & Exact \\
\hline
$K_0$ & $X_0 Z_1$ & 0.8374 & 1 \\
$K_1$ & $Z_0 X_1 Z_2$ & 0.8965 & 1 \\
$K_2$ & $Z_1 X_2 Z_3$ & 0.9136 & 1 \\
$K_3$ & $Z_2 X_3$ & 0.9697 & 1 \\
\end{tabular}
\end{center}

(\texttt{ibm\_kingston}, $4096$ shots per generator, 2026-07-29; job IDs
in the accompanying repository.) All four circuits were transpiled to the
same physical qubits used in the GHZ(3) runs of Section~5 (logical qubit
$i$ mapped to physical qubit $i$), including the same physical qubit $1$
whose published readout error that day was $23.4\%$. Of the four
generators, $K_0$, $K_1$, and $K_2$ all include physical qubit $1$ in
their support; $K_3$ is the only one that does not --- and $K_3$ is, by a
wide margin, the closest to the exact value. This is a second,
independent piece of evidence for the same conclusion Section~5 reached
by a different route (bit-flip statistics rather than stabilizer
expectation values): a single miscalibrated physical qubit, not the
number of qubits or gates in a generator, is the dominant factor in the
observed deviation from the exact prediction.

\section{Scope and limitations}

This note does not claim any new theorem. The entropy formula is due to
Fattal et al.\ (2004); the code-distance procedure is standard
(Gottesman, 1997). What is reported here is: (a) an exact,
dependency-light, floating-point-free implementation of both, (b)
explicit verification against an independent computation before any
result was trusted, (c) one small original observation (the
saturation scaling law of Section 4), offered honestly as a numerical
finding about a specific circuit model, with an unproven heuristic
explanation, not a general result about entanglement growth in quantum
circuits, and (d) two small, explicitly dated real-hardware comparisons
(Sections~\ref{sec:hardware}--\ref{sec:cluster-hardware}) that are
procedural and illustrative, not a noise-characterisation result --
though the two independently point to the same qualitative conclusion,
that individual qubit miscalibration dominates over circuit size in the
regime we tested.

\begin{thebibliography}{9}

\bibitem{Fattal2004}
D.~Fattal, T.~S.~Cubitt, Y.~Yamamoto, S.~Bravyi, and I.~L.~Chuang,
Entanglement in the stabilizer formalism,
\textit{arXiv:quant-ph/0406168}, 2004.

\bibitem{Gottesman1997}
D.~Gottesman,
Stabilizer Codes and Quantum Error Correction,
Ph.D.\ thesis, California Institute of Technology, 1997.
\textit{arXiv:quant-ph/9705052}.

\bibitem{NielsenChuang}
M.~A.~Nielsen and I.~L.~Chuang,
\textit{Quantum Computation and Quantum Information},
Cambridge University Press, 2000. Chapter 10.

\end{thebibliography}

\end{document}
